
A controlled comparison of token entropy, semantic entropy, and SelfCheckGPT-BERTScore on HaluEval-QA with LLaMA-2-7B-chat shows that all three methods fail to discriminate hallucinated from factual responses: AUROC values are 0.5000, 0.4829, and 0.3562 respectively. The failure of semantic entropy is traceable to the NLI clustering mechanism: microsoft/deberta-large-mnli assigns all five stochastic responses to distinct semantic clusters in 72.8\% of examples (aggregation rate = 0.272, 95\% CI [0.253, 0.292]), collapsing semantic entropy to a constant signal. SelfCheckGPT-BERTScore achieves below-random discrimination, consistent with — but not proven to be caused by — a label polarity inversion in HaluEval-QA's ChatGPT-generated confident hallucinations.

The cross-model comparison (Mistral-7B-Instruct, H-M3) was not executed, and cross-model generalizability of these findings is unknown.

The primary empirical contributions are: (1) first measurement of the NLI aggregation rate of deberta-large-mnli on HaluEval-QA (0.272, 95\% CI [0.253, 0.292]), establishing that non-degeneracy of semantic entropy cannot be assumed on short factual QA benchmarks; (2) documentation of below-random SelfCheckGPT-BERTScore discrimination (AUROC = 0.3562) on HaluEval-QA with LLaMA-2-7B-chat, with a proposed unverified explanation; (3) a matched-condition evaluation framework that makes all three UQ methods directly comparable on the same dataset and model.

Three experiments would directly extend these findings: substituting a QA-specific NLI model for deberta-large-mnli in the semantic entropy pipeline (zero new LLM inference required, using existing H-E1 stochastic samples); verifying the SelfCheckGPT polarity inversion by stratifying existing H-E1 consistency scores by label (zero new inference); and executing H-M3 with Mistral-7B-Instruct using the proven H-E1 codebase.

